\documentclass[11pt]{article}
\usepackage[margin=1in]{geometry}
\usepackage{amsmath,amssymb,amsthm,mathtools}
\usepackage[T1]{fontenc}
\usepackage{lmodern}
\usepackage[hidelinks]{hyperref}
\usepackage{microtype}
\newtheorem{theorem}{Theorem}[section]
\newtheorem{proposition}[theorem]{Proposition}
\newtheorem{corollary}[theorem]{Corollary}
\newtheorem{lemma}[theorem]{Lemma}
\theoremstyle{remark}
\newtheorem{remark}[theorem]{Remark}
\newcommand{\N}{\mathbb N}
\newcommand{\Q}{\mathbb Q}
\newcommand{\B}{\mathcal B}
\newcommand{\By}{\mathcal U}
\newcommand{\pred}{\operatorname{pred}}
\newcommand{\succs}{\operatorname{succ}}
\title{A singleton-supported cycle correction\\to Byrnes's differential-poset bound}
\author{Alec Kriebel\\
\small Independent researcher\\
\small ORCID: \href{https://orcid.org/0009-0001-9320-500X}{0009-0001-9320-500X}}
\date{Version 1.0.0 --- 22 October 2026\\\small Unrefereed preprint}
\begin{document}
\maketitle
\begin{abstract}
For every ordinary unweighted $r$-differential poset and $n\ge2$, we subtract
an explicit nonnegative term from Byrnes's cumulative upper bound for the
$n$th rank size. The term is the cycle rank of the preceding incidence
subgraph supported on elements with a singleton successor. Singleton
witnesses transfer these cycles into the next incidence graph by replacing
stars with trees having the same terminals. The correction can be positive
and attained on standard reflected Young-prefix examples, and its coefficient
one is sharp. A singleton-successor condition gives the conjectured local
recurrence; two such successors give a conditional equality characterization.
The result answers the printed request for an improvement of Byrnes's
inequality in an incidence-dependent form. It does not prove the general
local recurrence, improve the global Fibonacci envelope, or give a
combinatorial proof of strict rank growth.
\end{abstract}

\section{The improvement question and attribution}
Majumder's problem in Oberwolfach Report 23/2018, printed p.~1452
\cite{Majumder}, asks whether Byrnes's bound
\begin{equation}\label{eq:byrnes}
 p_n\le \By_n:=1+r\sum_{k=0}^{n-2}p_k+(r-1)p_{n-1}
\end{equation}
can be improved toward Stanley's proposed local inequality
\begin{equation}\label{eq:local}
 p_n\le r p_{n-1}+p_{n-2}.
\end{equation}
The same contribution separately asks for a combinatorial strict-growth
proof. Our claim concerns the improvement request only. The correction uses
incidences of the preceding ranks; no rank-number-only restriction is imposed
in the printed request.

The operator framework, common-cover restrictions and reflection construction
are classical \cite{Stanley}. Byrnes proves \eqref{eq:byrnes} on the way to the
global Fibonacci envelope and states \eqref{eq:local} separately
\cite[Chapter~5]{Byrnes}. Incidence and hypergraph dimension sums already
appear in \cite[Remark~8]{SZ}; the cycle identities below repackage these
counts. Cycle-space monotonicity and terminal-preserving tree replacement
are elementary graph facts. The substantive assertion proved here is the
simultaneous transfer of the singleton-supported cycle rank, not a new
general graph method.

\newpage
\section{Incidence graphs and exact slack}
An ordinary $r$-differential poset is a locally finite $\N$-graded poset $P$
with a unique minimum $\widehat0$, finite ranks $P_j$, and ordinary unweighted
up and down operators over $\Q$ satisfying $DU-UD=rI$, where $r$ is a positive
integer. Write $p_j=|P_j|$, and let $d(x)$ and $u(x)$ count the lower and upper
covers of $x$. Equivalently, the integer incidence counts satisfy
\begin{equation}\label{eq:counts}
 u(x)=d(x)+r,\qquad
 |\pred(x)\cap\pred(y)|=|\succs(x)\cap\succs(y)|
\end{equation}
for distinct elements in a common rank. We do not interpret this identity
only modulo a characteristic. The terminal equality to one in the OWR
common-cover axiom is a typographical error: zero common covers are allowed
under the standard definition \cite{Stanley,GV}.

\begin{lemma}\label{lem:basic}
Distinct same-rank elements have at most one common lower cover and at most
one common upper cover. Every consecutive-rank incidence graph is connected.
\end{lemma}
\begin{proof}
If a pair has two common upper covers, \eqref{eq:counts} gives two common
lower covers. This lower pair has the original pair as common upper covers.
Iterating down the ranks contradicts the unique rank-zero element. The
common lower count is equal to the upper count.

Let $\B_j$ be the bipartite incidence graph on $P_{j-1}\sqcup P_j$.
The graph $\B_1$ is a star with $r$ leaves. If $\B_j$ is connected, its
projection on $P_j$, joining elements with a common predecessor, is connected
(including the one-vertex case). By \eqref{eq:counts} each such adjacency has
a common successor, so all $P_j$ vertices lie in one component of $\B_{j+1}$.
Every vertex of $P_{j+1}$ has a lower cover, proving the induction.
\end{proof}

For a finite graph $H$, retaining isolated vertices, define
$\beta(H)=|E(H)|-|V(H)|+c(H)$, where $c(H)$ is its component count.
This is its cycle-space dimension over any fixed field; it is nonnegative
and monotone under subgraph inclusion. Put $b_j=\beta(\B_j)$ for $j\ge1$.

\begin{proposition}\label{prop:slack}
For $j\ge1$, the number of edges of $\B_j$ is
$e_j=r\sum_{k=0}^{j-1}p_k$, and
\begin{align}
 b_j&=r\sum_{k=0}^{j-1}p_k-p_{j-1}-p_j+1,\label{eq:bj}\\
 p_n&=\By_n-b_n &&(n\ge2),\label{eq:slack}\\
 b_n-b_{n-1}&=r p_{n-1}+p_{n-2}-p_n &&(n\ge2).\label{eq:increment}
\end{align}
\end{proposition}
\begin{proof}
Summing the degree identity gives $e_j=e_{j-1}+rp_{j-1}$, with $e_0=0$.
Connectivity gives \eqref{eq:bj}; rearranging and subtracting gives the
remaining identities.
\end{proof}
Thus Byrnes's bound is $b_n\ge0$, whereas \eqref{eq:local} requires
$b_n\ge b_{n-1}$. Nonnegative cycle rank alone does not give monotonicity.

\newpage
\section{The supported-cycle theorem}
Fix $n\ge2$ and set $Z=P_{n-2}$, $X=P_{n-1}$. Define
\[
 S=\{z\in Z:\text{some }a\in X\text{ has }\pred(a)=\{z\}\}.
\]
Such an $a$ is a \emph{singleton successor} of $z$; it does not mean that
$z$ has only one successor. Let $H_S$ retain all vertices of $X$, precisely
$S$ on the other side of $\B_{n-1}$, and all corresponding incidence edges.
Set
\begin{equation}\label{eq:q}
 q_n=\beta(H_S)
 =\sum_{z\in S}(u(z)-1)-p_{n-1}+c(H_S).
\end{equation}

\begin{theorem}\label{thm:main}
For every ordinary $r$-differential poset and every $n\ge2$,
\begin{equation}\label{eq:transfer}
 0\le q_n\le b_{n-1},\qquad b_n\ge q_n.
\end{equation}
Consequently
\begin{align}
 p_n&\le1+r\sum_{k=0}^{n-2}p_k+(r-1)p_{n-1}-q_n,\label{eq:refined}\\
 p_n&\le r p_{n-1}+p_{n-2}+b_{n-1}-q_n.\label{eq:residual}
\end{align}
The correction is determined by the two preceding ranks, before $P_n$ is
constructed, with no further hypothesis on the poset.
\end{theorem}
\begin{proof}
The first two inequalities follow from $H_S\subseteq\B_{n-1}$.
For each $z\in S$, choose $a_z\in X$ with predecessor set $\{z\}$ and put
$N_z=\succs(z)\subseteq X$. For every $y\in N_z\setminus\{a_z\}$,
the pair $a_z,y$ has unique common predecessor $z$, hence unique common
successor $t\in P_n$. Every lower cover $w$ of $t$ lies in $N_z$:
if $w\ne a_z$, its common successor with $a_z$ forces a common predecessor,
which must be $z$.

Take every nonsingleton upper vertex $t$ covering $a_z$, with all its
incidence edges, and all terminals $N_z$. Two different such $t$ cannot
share another terminal, by Lemma~\ref{lem:basic}. Each terminal other than
$a_z$ occurs exactly once. Thus these stars, joined at $a_z$, form a tree
$T_z$ connecting precisely $N_z$. If $|N_z|=1$, take the one-vertex tree.
A rank-$n$ internal vertex cannot belong to both $T_z$ and $T_w$ for $z\ne w$:
it would cover the witnesses $a_z,a_w$, whose predecessor sets are disjoint,
contrary to \eqref{eq:counts}. Different trees can share only $X$ terminals.

Let $K$ be their union together with every vertex of $X$, including isolates.
Replace each star of $H_S$ centered at $z$ by $T_z$.
If $T_z$ has $a$ internal vertices, it has $|N_z|+a-1$ edges.
The replacement changes both vertex and edge counts by $a-1$ and preserves
the connections among terminals. Internal vertices from different
replacements are disjoint, so components and cycle rank are preserved:
$\beta(K)=\beta(H_S)=q_n$. Since $K\subseteq\B_n$, cycle-space inclusion gives
$b_n\ge q_n$. Equations~\eqref{eq:slack} and \eqref{eq:bj} now prove the bounds.
\end{proof}

\newpage
\section{When the local bound and equality follow}
\begin{corollary}\label{cor:all}
If every $z\in P_{n-2}$ has a singleton successor in $P_{n-1}$, then
$p_n\le r p_{n-1}+p_{n-2}$. In particular this holds for every admissible next
extension whenever rank $n-1$ was formed by reflection-extension.
\end{corollary}
\begin{proof}
Here $S=Z$, so $q_n=b_{n-1}$. Reflection creating rank $n-1$ adds $r$
singleton successors of every rank-$(n-2)$ element; see the construction in
Section~\ref{sec:sharp}. The extension to rank $n$ may be arbitrary.
\end{proof}

\begin{theorem}\label{thm:equality}
Suppose $n\ge3$ and every $z\in P_{n-2}$ has at least two distinct singleton
successors in $P_{n-1}$. Equality in the local bound holds if and only if the
extension to rank $n$ is reflection-extension, up to relabeling that rank.
\end{theorem}
\begin{proof}
Choose one witness per $z$ in the main proof. Now $K$ is connected and contains
all $X$ vertices. Each unused upper vertex $t$ attaches by its complete star
to $X$, adding $d(t)-1$ to cycle rank. Equality means
$b_n=b_{n-1}=\beta(K)$, so every unused upper vertex is singleton.

For a fixed $z$, take distinct singleton successors $a=a_z,b$.
Every nonsingleton upper vertex covering $b$ belongs to $K$ and specifically
to $T_z$: if it belonged to $T_w$, its common predecessor with the witness
$a_w$ would force $w=z$. Every internal vertex of $T_z$ covers $a$.
Let $h$ be the unique common successor of $a,b$.
For each $y\in N_z\setminus\{b\}$, its common successor with $b$ is
nonsingleton, lies in $T_z$, and covers $a$. Uniqueness for $a,b$ forces it
to be $h$. Hence $h$ covers all $N_z$, and no second nonsingleton internal
vertex can share a pair of terminals.

Since $n\ge3$, each $z$ has a predecessor and $u(z)\ge r+1\ge2$.
There is therefore precisely one nonsingleton vertex $h_z$ over each $N_z$.
They are distinct by the witness argument. Each $x\in X$ belongs to exactly
$d(x)$ of these neighborhoods; the degree identity leaves exactly $r$
singleton successors over $x$. This is reflection-extension.
Conversely reflection has $|Z|+r|X|$ upper vertices and gives equality.
\end{proof}

The following established lemma is \cite[Proposition~3.4]{GV}; its short
counting proof is included to make the consequence self-contained.
\begin{lemma}[Gaetz--Venkataramana]\label{lem:rankone}
Every rank-one element has at least two singleton successors.
\end{lemma}
\begin{proof}
There are $r$ rank-one elements, each with $r+1$ successors. At a fixed $z$,
nonsingleton successors use disjoint nonempty subsets of the other $r-1$
rank-one elements; otherwise a pair would have two common successors.
There are at most $r-1$ nonsingleton successors, leaving at least two.
\end{proof}
Thus, for every $r\ge1$, $p_3\le rp_2+r$, with equality precisely for reflection
from rank two. This is a consequence of the transfer/equality theorem and
the credited lemma, not a claim that the lemma is new.

\newpage
\section{Standard examples giving strictness and sharpness}\label{sec:sharp}
We recall Wagner reflection, as presented in \cite[Proposition~6.1 and the
following paragraph]{Stanley}. Given a valid prefix through rank $m\ge1$,
add one vertex over the entire successor set in $P_m$ of each element of
$P_{m-1}$, and add $r$ distinct singleton successors over each element of
$P_m$. Each current $x$ gets $d(x)+r$ successors; a current pair gets one
common new successor for each common predecessor. Every new vertex has
a predecessor and the rank is finite. Thus the prefix prolongs indefinitely
to an ordinary differential poset. Its next rank has $p_{m-1}+rp_m$ elements.
The constructions below use this classical operation and classical Young
prefixes, rather than new poset families.

For $r=3$, take rank one $\{a,b,c\}$. Put one rank-two element over each of
$\{a,b\},\{a,c\},\{b,c\}$ and two singleton successors over each rank-one
element. Each rank-one element has four successors and every pair has one.
This is the prefix of $Y^3$ through rank two, with ranks $(1,3,9)$.
Reflect to rank three and continue reflecting. Then
\[
 (p_0,p_1,p_2,p_3)=(1,3,9,30),\qquad
 e_2=12,\quad e_3=39,\quad b_2=b_3=q_3=1.
\]
Byrnes's rank-three bound is $31$ and \eqref{eq:refined} gives the attained
bound $30$.

For $r=1$, use Young's lattice $Y$ through rank five and reflect twice,
then continue reflecting. Young's lattice consists of integer partitions,
ordered by addition of one box. Its first six rank sizes are
$(1,1,2,3,5,7)$, so the extended ranks are
\[
 (1,1,2,3,5,7,12,19).
\]
Here $e_6=\sum_{k=0}^5p_k=19$ and $p_5+p_6=19$, hence $b_6=1$.
Reflection at rank six gives a singleton successor for each rank-five element,
so $q_7=b_6=1$; reflection at rank seven gives $b_7=1$.
Byrnes's bound $20$ is reduced to the attained bound $19$.
The standard reflection argument proves infinite realization in both cases.

In either example $p_n=\By_n-q_n$ with $q_n=1$. Replacing the coefficient
one in \eqref{eq:refined} by any $c>1$ would yield a false bound
$p_n\le\By_n-cq_n<p_n$. This proves coefficient sharpness, not optimality
of every possible use of the preceding incidence data.

\subsection*{Boundary cases and the exact remaining gap}
No $b_0$ is needed: $b_1=0$. At $n=2$, $Z=\{\widehat0\}$ and every rank-one
element is a singleton successor, so $q_2=0$ and $p_2\le r^2+1$.
If also $r=1$, the replacement tree has one terminal and no internal vertex.
For empty $S$, $H_S$ is the isolated vertex set $X$ and $q_n=0$.

Always $\By_n-(rp_{n-1}+p_{n-2})=b_{n-1}$, so the refined right side lies
between Byrnes's and the proposed local right sides. The unremoved allowance
is exactly $b_{n-1}-q_n$. It can be positive: Young's lattice at $n=6$ has
$b_5=1,q_6=0$, refined bound $13$, local bound $12$, and actual $p_6=11$.
The general recurrence still requires $b_n\ge b_{n-1}$. No combinatorial
strict-growth argument follows here; the standard strict-growth statement
also has the exception $r=n=1$, when $p_0=p_1=1$.

\newpage
\section{Verification, related work, and research status}
The companion archive contains the complete \LaTeX\ source, all eight
unchanged submitted artifacts, the independently written standard-library
exact checker, complete expected outputs, a portable offline reproduction
runner, source/attribution records and content manifests. No third-party
full text, ledger data or private coordination is included.

The checker reconstructs predecessor sets, verifies every integer commutator,
connectivity, exact cycle identities and the actual witness-tree replacements.
Its edge-clique partition enumeration gives labeled prefix counts
$1,1,1,1,2,7,60$ for $r=1$ through rank seven, $1,1,4$ for $r=2$ through rank
three, and $1,2,288$ for $r=3$ through rank three. These are not isomorphism
counts. It checks reflected prefixes for $1\le r\le5$ through rank five,
Young's lattice through rank twelve, and both sharp examples.
The 5764 witness selections use first and last choices in every reflection-only
run, irrespective of product size. Enumerated-prefix, Young-lattice and separate
sharp-example runs exhaust products of at most 2000 choices and use first/last
above 2000. The proof establishes arbitrary witness choice.

Negative controls reject deletion of a top vertex and a false minimum degree,
and refute subtracting $2q_n$ in the sharp examples. The relaxed layer sizes
$(1,3,3,7)$ satisfy the two middle commutators for $r=1$ but violate the
rank-zero commutator; they do not form a differential-poset prefix.
Their pair-block projection has a triangle and their bipartite incidence
graph a six-cycle. Normal, \texttt{-O}, and \texttt{-OO} runs agree byte for
byte. These finite computations support implementation and example
correctness; the universal theorem and infinite realization rest on the
written proofs.

Singleton and crown techniques have earlier antecedents
\cite{Lewis,MR}; the exact rank-one two-singleton lemma used here is
\cite[Proposition~3.4]{GV}. Reflection of arbitrary Young prefixes is already
covered by \cite{Stanley}. Related upper-extremal investigations include
Qing's 2008 thesis and Wright's 2019 slides/abstract \cite{Qing,Wright};
their global-envelope or restricted-prefix scope is distinct from the
supported-cycle assertion. Byrnes also uses cycle arguments under other
hypotheses, so this note asserts no new general use of cycles in the subject.

A dated bounded primary-source comparison, completed 22 October 2026,
found no earlier equivalent supported-cycle upper-bound theorem in the
inspected material. Historical firstness, first application and present
worldwide openness are not certified. Source versions, inspected portions
and access limits are recorded in the companion attribution dossier.
The final publisher bodies of Gaetz--Venkataramana and Miller were not
obtained; their versioned primary preprints were inspected. No subsequent
full proof corresponding to Wright's talk was located, and subscription
indexing and unpublished work remain coverage limits. The classification
of this result concerns the literal printed improvement request, not every
stronger rank question.

\subsection*{AI involvement, review, and licenses}
AI tools were extensively involved in research planning, mathematical
exploration and proof development/reconstruction, code and exact computations,
literature and source investigation, writing, editing, and adversarial
mathematical and priority auditing. Alec Kriebel is the named human author.
This manuscript is unrefereed and has not undergone professional independent
human peer review, journal acceptance or formal proof-assistant certification.
Internal AI reviews do not supply those certifications. A fresh whole-package
review is a separate release gate. Original prose is offered under CC BY 4.0;
the companion code is MIT licensed. Cited works retain their respective rights.

\newpage
\begin{thebibliography}{99}\small\raggedright
\bibitem{Stanley} R.~P. Stanley, Differential posets, \emph{J. Amer. Math.
Soc.} \textbf{1} (1988), 919--961.
\href{https://doi.org/10.1090/S0894-0347-1988-0941434-9}{doi:10.1090/S0894-0347-1988-0941434-9}.
\bibitem{Byrnes} P. Byrnes, \emph{Structural Aspects of Differential Posets},
Ph.D. thesis, University of Minnesota, 2012, Chapter~5.
\url{https://conservancy.umn.edu/server/api/core/bitstreams/45a1bae8-402f-42a7-ad67-4caa7b55d104/content}.
\bibitem{Majumder} P. Majumder, Rank sizes of Differential Posets, in
Open Problem Session, \emph{Oberwolfach Reports} 23/2018, printed pp.~1452--1453.
Report doi:\href{https://doi.org/10.4171/OWR/2018/23}{10.4171/OWR/2018/23}.
\url{https://ems.press/content/serial-article-files/46745}.
\bibitem{SZ} R.~P. Stanley and F. Zanello, On the rank function of a
differential poset, \emph{Electron. J. Combin.} \textbf{19}(2) (2012), P13.
\href{https://doi.org/10.37236/2258}{doi:10.37236/2258}.
\bibitem{GV} C. Gaetz and P. Venkataramana, Path counting and rank gaps in
differential posets, \emph{Order} \textbf{37} (2020), 279--286 (online 2019).
\href{https://doi.org/10.1007/s11083-019-09504-4}{doi:10.1007/s11083-019-09504-4}.
Inspected primary preprint: \url{https://arxiv.org/abs/1806.03509v3}.
\bibitem{Lewis} J.~B. Lewis, \emph{On Differential Posets}, undergraduate thesis,
Harvard University, 2007, Section~3.
\url{https://bpb-us-e1.wpmucdn.com/blogs.gwu.edu/dist/5/693/files/2018/04/JBLHarvardSeniorThesis-20drube.pdf}.
\bibitem{MR} A.~R. Miller and V. Reiner, Differential posets and Smith normal
forms, \emph{Order} \textbf{26} (2009), 197--228.
Inspected preprint: \url{https://arxiv.org/abs/0811.1983v1}.
\bibitem{Qing} Y. Qing, \emph{Differential Posets and Dual Graded Graphs}, S.M. thesis, Massachusetts
Institute of Technology, September 2008, Chapters~2--3.
\url{https://web.math.utk.edu/~yqing/MasterThesis.pdf}.
\bibitem{Wright} B. Wright, \emph{Maximal Extensions of Differential Posets}, MIT PRIMES conference
slides, 19 May 2019, and conference abstract, p.~12.
\url{https://math.mit.edu/research/highschool/primes/materials/2019/conf/9-1-Wright.pdf}.
\bibitem{Miller} A.~R. Miller, Differential posets have strict rank growth:
a conjecture of Stanley, \emph{Order} \textbf{30} (2013), 657--662.
Inspected preprint: \url{https://arxiv.org/abs/1202.3006v1}.
\end{thebibliography}
\end{document}
